\documentclass{article}
\usepackage{graphicx} % Required for inserting images

\title{lec427}
\author{Dylan Tang}
\date{April 2026}

\begin{document}

\maketitle

\section{Quiz Format}
\begin{itemize}
    \item Mix of coding and theory
    \item \texttt{scikit fit, train\_test\_split}
\end{itemize}

\section{Autoregression}

\subsection{Case Study: Linear Regression on Month Number and Sunspots}
\begin{verbatim}
    from sklearn.linear_model import LinearRegression
    from sklearn.model_selection import train_test_split
    from sklearn.metrics import r2_score
    import numpy as np

    X = df[['Month_No']].to_numpy()
    Y = df[['Sunspots]].to_numpy()

    Xtrain, Xtest, Ytrain, Ytest = train_test_split(X,Y)

    reg = LinearRegression()
    reg.fit(Xtrain, Ytrain) # train linear model in training set

    Yhat = reg.predict(Xtest) # predict on test set

    result = {'r2_score': r2_score(Ytest,Yhat)}
    print(result)
\end{verbatim}
\begin{itemize}
    \item Returns a very low $R^2$ score, so month is not a good predictor of sunspots.
    \item But, sunspots on day $i$ is a good predictor of sunspots on day $i+1$
\end{itemize}

\newpage

\subsection{AR(1) Model}
\begin{verbatim}
    from sklearn.linear_model import LinearRegression
    from sklearn.model_selection import train_test_split
    from sklearn.metrics import r2_score

    df['Sunspots_prev'] = df['Sunspots].shift(1)
    df = df.dropna(how='any') # Drop the first value

    X = df[['Sunspots_prev']].to_numpy()
    Y = df[['Sunspots']].to_numpy()

    Xtrain, Xtest, Ytrain, Ytest = train_test_split(X,Y)

    reg = LinearRegression()
    reg.fit(Xtrain, Ytrain) # train linear model in training set

    Yhat = reg.predict(Xtest) # predict on test set

    result = {'r2_score': r2_score(Ytest,Yhat)}
    print(result)
\end{verbatim}
\begin{itemize}
    \item When regress texttt{Sunspots} on texttt{Sunspots\_prev}, the $R^2$ improves to 0.84
\end{itemize}

\begin{itemize}
    \item However, standard \texttt{train\_test\_split} is not enough, as \texttt{train\_test\_split} often randomly shuffles rows, which can accidentally train on the future and test on the past.
\end{itemize}
\subsection{AR(1) model with Time-series Train-Test Split}
\begin{verbatim}
    from sklearn.linear_model import LinearRegression
    from sklearn.model_selection import train_test_split
    from sklearn.metrics import r2_score
    from sklearn.model_selection import TimeSeriesSplit
    
    df['Sunspots_prev'] = df['Sunspots].shift(1)
    df = df.dropna(how='any') # Drop the first value

    X = df[['Sunspots_prev']].to_numpy()
    Y = df[['Sunspots']].to_numpy()

    tscv = TimeSeriesSplit()
    for train_index, test_index in tscv.split(X):
        Xtrain, Xtest = X[train_index], X[test_index]
        Ytrain, Ytest = y[train_index], y[test_index]

    reg = LinearRegression()
    reg.fit(Xtrain, Ytrain) # train linear model in training set

    Yhat = reg.predict(Xtest) # predict on test set

    result = {'r2_score': r2_score(Ytest,Yhat)}
    print(result)
\end{verbatim}

\section{Longitudinal vs Cross-sectional Models}
\begin{itemize}
    \item \textbf{Longitudinal model (L model)} - uses $y_{i-1}$ to predict $y_i$
    \item \textbf{Cross-sectional model (C model)} - use predictors $x_{i}$s to predict $y_i$
\end{itemize}

\subsection{Capital Asset Pricing Model (CAPM)}
\begin{itemize}
    \item Predicting market prices can be broken down into a C model (uses outside predictors) and the L model (uses market prices from the previous day)
    \item The L model is often taken as given, and building the C model to predict where the L model fails to predict is the main point of interest
\end{itemize}

\subsection{texttt{scikit} One Hot Encoder}
\begin{itemize}
    \item Turns non-numerical predictor levels into numerical values
\end{itemize}
\begin{verbatim}
    from sklearn.preprocessing import OneHotEncoder

    feat1 = OneHotEncoder(sparse=False) #
    Xsmoker = feat1.fit_transform(df["smoker"].to_numpy().reshape(-1,1)) # one hot encode the "smoker" variable, and reshape into column
    X = np.hstack([X, Xsmoker]) # append column to X

    ...
\end{verbatim}
\end{document}